\documentclass[12pt]{article}

\usepackage[a4paper, total={6in, 8in}]{geometry}
\usepackage{textcomp}

\begin{document}
\setlength{\parindent}{0pt}

\def \t {\textrightarrow}
\def \v {\vspace{1em}}
\def \bi {\begin{itemize}}
\def \ei {\end{itemize}}
\def \s[#1] {\section*{#1}}
\def \ss[#1] {\subsection*{#1}}
\def \sss[#1] {\subsubsection*{#1}}

\s[Guerra fredda]
Programma della nuova frontiera \t il 22 Novembre 1963 Kennedy viene ucciso, probabilmente perchè la sua politica contro discriminazioni non piaceva \\
Infatti anche Martin Luther King viene ucciso nel 68, e poi anche Robert Kennedy (il fratello) che nel 68 stava per diventare il presidente degli USA \\
Si vede anche una leggera apertura della Chiesa \t uno dei protagonisti della distensione è Papa ? \t nel 63 pubblica enciclica "Pacecm in ternis" \t richiama a risoluzioni questioni in modo diplomatico \\
Alla morte di Kennedy diventa presidente Johnson

\ss[Guerra in Vietnam]
Vietnam aveva lottato per la sua indipendenza gia da prima \t contro i francesi, che abbandonano l'indocina nel 64 con accordi di Ginevra \t e Vietnam diviso in due :
\bi
\item Ocimin guida Nord, comunista
\item filoccidentale al sud, governata da un cattolico
\ei
Il cattolico è pero ha caratteri autoritari, e non piace quindi a popolazione per lo piu buddista \t quindi si crea il movimento dei Vietcong che si oppone al governo filoccidentale \\
I Vietcong, movimento anche di guerriglia, vengono supportati dal Vietnam del Nord \\
Ora gli USA temono che tutta quella zona possa diventare comunista \t ma era importante mantenere una parte dell'Asia \\
Nel 56 quindi USA decidono di intervenire in Vietnam \t e manadano i consiglieri militari, che sono soldati in realtà \t vengono mandati per supportare il Sud \\
Il numero dei consiglieri verrà incrementato, e arriveranno ad essere 30 mila uomini, che sono in Vietnam \\
Golfo del Tonchino \t si verifica incidente, in cui la marina vietnamita attacca (o forse si sono scontrari) un carpadiniere americano???? \\
Questo incidente è stata la svolta che ha permesso agli USA di aumentare la usa presenza \t dal 56 in poi, USA aveva cercato ogni scusa in realta \t c'è chi dice che sia stato graduale, chi che questa è stata la svolta

\v

Nel 1965 Johnson decide di passare a un azione di forza \t quindi bombarda il Nord senza dichiarazione di guerra \\
D'ora in poi USA continuano a mandare militari, e arriveranno ad essere mezzo milione \t ma non vincono, e si scontrano con Nord e anche Vietcong al sud \\
Inoltre anche URSS e Cina supportavano movimenti comunisti \\
Nel 68 i Vietcong lanciano l'offesiva del Tet (capodanno vietnamita) verso il sud \t non ha molti guadagni, ma si capisce che guerra sara lunga e non è detto che USA vincera \\
L'opinione pubblica americana detesta questa guerra \t quindi nel 68 i bombardamenti vengono fermati, e Johnson dichiara di non candidarsi per presidenziali \\
Il suo successore, Nixon, avviera le trattative di pace e ridurra l'impegno militare americano \t la linea chepreval è quella di ritirata, e di lasciare il Vietnam nella sua situazione \t non si vuole far credere che USA non ce la fa \\
Nixon richiama esercito, e cerchera di isolare i Vietcong coinvolgendo nel conflitto Laos e Cambogia \t si ritira, pero cerca di supportare il Vietnam del Sud e rafforzarne l'esercito \t ma fallira \\
Nel 73 firmato armistizio a Parigi, che sancisce ritiro del contingente americano \t non è una pace, me un "me ne vado" \\
Nel 75 ? ci sara riunficazione del Vietnam \t perche Vietcong e Nord entrano a Saibon (capitale del Sud)

\v

La guerra del Vietnam è uno dei momenti + difficili della storia recente USA \t è l'unica guerra che hanno perso \\
Il loro intervento viene condannato a livello internazionale, e USA devono affrontare contestazione interna importante \t l'opinione pubblica si chiede perche giovani americani dovevano andare a combattere in Vietnam \\
Inoltre c'era anche domanda, del perche piu grande potenza al mondo non riusciva a vincere in un paese sottosviluppato \\
Governo dice che guerra era necessaria per combattere comunismo e impedire espansione in tutta asia di comunismo \\
Inoltre dice che non era guerra tradizionale, ma di guerriglia \t le truppe americane non sono pronte a questo tipo di combattimento, e non conoscono territorio \t governo si difende cosi \\
Comunque grande opposizione dell'opinione pubblica

\v

Americani si ritirano, e anche in Laos e Cambogia nascono regimi comunisti \t in Cambogia la Cina interviene, che teme influenza troppo sovietica \\
Sostiene quindi Polpot, leader dei Qmer Rossi \t guerriglieri comunisti \\
La cina cerca di bilanciare influenza sovietica nella zona, che influenza ??? \\
Qmer creano repressione fortissima, quasi di sterminio \t interviene Vietnam per allontanare Qmer a instaurare regime fedele a quello Vietnamita (nel 78) \\
Nel 79 la Cina attacca il Vietnam \t pero senza successo, perchè la Cambogia rimane nell'orbita vietnamita

\ss[La Cina]
In Cina nel 49 c'è stata la rivoluaizone, Mao al potere crea repubblica popolare cinese \\
Cina da inizio a forte industrializzazione \t con supporto anche di tecnici dall'URSS \t senza pero decollo industriale \\
Anche attuata collettivizzazione delle terre \t nel 1950 vara riforma agraria, migliore di quella di stalin \t che distribuisce la guerra ai contadini \\
Nascono aziende agricole che si uniscono in cooperative controllate dallo stato \t ma senza decollo del'agricoltura \\
I cinesi vedono anche un'impressionante crescita demografica \t nel 54 la popolazione cresce di 100 milioni \t questo impedisce anche decollo agricoltura, perchè cio che viene prodotto deve sfamare molti \\
Politica del Grande Balzo in Avanti \t avviata da Mao \t le cooperative dovevano diventare autosufficienti in tutti i settori, e quindi diventare + efficienti \\
Propaganda martellante, e controlli molto rigidi su queste cooperative, che dovevano essere redditizie

\v

Fallimento totale \t 30 milioni di morti di fame \\
Inotlre Cina voleva prendere distanze da URSS nonostante iniziale collaborazione \t cosi l'urss decide di non fornire assistenza alla cina in ambito nucelare, vuole mantenere supremazia \\
Nel 64 la cina quindi si dota della bomba atomica in autonomia, in contesto di tensione con urss \t e urss cerca di muovere movimento comunista internazionale contro la cina, senza successo \\
Nasceranno discussioni sui confini anche tra cina e urss \t la tensione tocca tutti gli ambiti \t è un conflitto egemonico e ideologico (comunismo industriale contro agricolo)

\v

Fallisce politica di Mao, e quindi diventano + popolari i moderati \t Mao reprime fortemente, e culmina con la rivoluazione culturale = tentativo dei riformisti di contestare ogni forma di autorità precostituita \\
Protesta viene repressa con un anti rivoluzione culturale \t e incita i giovani a opporsi a trasformazione del comunismo in capitalismo (cosi viene venduto ) \\
Crea anche campi di rieducazione (tipo gulag) \t vittime sono 6 milioni \\
Mao voleva cambiare del tutto la cutlura del popolo cinese \t si voleva evitare che idea riformista si potesse diffondere, e bisognava andare invece verso modello comunista \\
Il simbolo di questa lotta è il "Libretto rosso" di Mao \t e l'attuamento culturale sara di ispirazione per molti gruppi giovanili

\v

Questo incide sull'economia e su coesione del gurppo dirigente \t e quindi Mao torna sui suoi passi \t elementi + radicali del partito messi ai margini, repressione delle Guardie Rosse vengono deposti \\
Questo verra fatto anche grazie al primo ministro, che genera dei cambiamenti non aggressivi \\
Cambiamento anche in politica estera \t cina si trova isolata a causa di ostilita urss, e non ha alleati \\
Quindi cina si apre agli USA \t cambiamento radicale \t nel 1972 Nixon va a Pechino, e Cina viene ammessa all'ONU \t si va verso una normalizzazione dei rapporti \\
Mao muore nel 76

\s[Movimento del 68]
Esplode nelle univeristà americane legata all'accademia \t ma poi anche sul vietnam \\
Nel 68 i giovani diventano un nuovo gruppo sociale \t diventa un gruppo sociale con dei caratteri e dei valori comuni, ed è un gruppo che si contrappone alla vecchia società \\
La svolta inizia gia negli anni 50 in realta \t quando si manifesta tra i giovani un ansia esistenziale senza precedenti \t e per la prima volta giovani diventano un gruppo sociale \\
Diventano anche una nuova categoria di consumatori \t USA ha sviluppo economico, e questo permette di generare posti di lavoro che si aprono ai giovani \t quindi diventano consumatori ?? \\
Cambia quindi al mercato a fine anni 50 \t magliette, scarpe da ginnastica, etc. \t nasce mercato orientato ai giovani \\
Il sistema viene contestato in quanto tale \t tutto cio che è preordinato, che è il fondamento della societa ma non è stato scelto \t viene messo in discussione \\
Nascono i miti di Cheguevara, Mancuse, etc. \\
Nasce l'idea che la cultura passi attraverso l'associazionismo e il confronto anche \\
Anche in europa, a Berlino e Parigi \t nasce senso di solidarietà tra tutti i giovani del mondo \t e allo stesso tempo accadono diverse eventi centrali

\v

Nel 79 in Iran c'è rivoluzione che scaccia lo Scha \t che era alleato degli USA, e Iran era fortemente occidentalizzato (imposizione occidentalista dello Scha) \\
Quindi aveva nemici, gruppi di sinistra e islamici sono contrari \t nel 79 rivoluzione, lui va in esilio e sale al potere Ayatollah Comeini \t esponente sciita, e nasce Repubblica Islamica \t è una teocrazia, e si ispira a un osservazione rigorosa del corano \\
I + moderati e progressisti vengono allontanati, e anche eliminati spesso \\
La repubblica islamica cerca di creare consenso e muove la popolazione contro il nemico esterno, ovvero l'USA \\
È una repubblica molto aggressiva, infatti a Teheran viene occupata l'ambasciata americana \t e diventano imprigionati \\
Lui muore nel 89 \t fino ad arrivare a Camenei, e ora non si sa \t ma da quel momento l'orientamento dell'Iran è fondamentalista \\
L'Iran di Comeini è diffidente anche verso URSS \t perchè l'URSS è diffidente nei confronti dell'Iran, e teme influenza dei musulmani comunisti \\
Iran appoggia anche guerriglia libanese, ma il nemico peggiore è l'Iraq \t ci sono contrasti politici e territoriali \\
Nell 80 c'è guerra feroce che dura fino all 88 \t si combatte principalmente sul confine, 1 milione di morti \t e non c'è vincitori \t tutti e due escono stremati e firmano armistizio
\end{document}
